\documentclass[journal]{IEEEtran}

\usepackage{amsmath, amsthm, amssymb, amsfonts}
\usepackage{algorithm}
\usepackage{algpseudocode}
\usepackage{graphicx}
\usepackage{booktabs}
\usepackage{cite}
\usepackage{hyperref}

\newtheorem{theorem}{Theorem}
\newtheorem{lemma}{Lemma}
\newtheorem{proposition}{Proposition}
\newtheorem{corollary}{Corollary}
\newtheorem{definition}{Definition}
\newtheorem{assumption}{Assumption}
\newtheorem{remark}{Remark}

\newcommand{\NN}{\mathcal{N}}
\newcommand{\HH}{\mathcal{H}}
\newcommand{\AA}{\mathcal{A}}
\newcommand{\MM}{\mathcal{M}}
\newcommand{\GG}{\mathcal{G}}
\newcommand{\VV}{\mathcal{V}}
\newcommand{\EE}{\mathcal{E}}
\newcommand{\PP}{\mathcal{P}}
\newcommand{\CC}{\mathcal{C}}
\newcommand{\Phii}{\Phi}

\title{Resilient Average Consensus with Adversaries: \\
       An Exposition and Open-Source Implementation}

\author{Alan~John%
\thanks{This document is an exposition adapted from L.\ Yuan and H.\ Ishii,
``Resilient Average Consensus with Adversaries via Distributed Detection and
Recovery,'' arXiv:2405.18752v1, May 2024.  The contributions of this exposition
are (i)~a self-contained re-derivation of the protocol with full pseudocode,
and (ii)~a paired open-source implementation in Python and MATLAB that
reproduces the paper's numerical examples.  The original algorithms,
theorems, and proofs are due to Yuan and Ishii.}}

\markboth{Resilient Average Consensus -- Exposition by Alan John, April 9, 2025}{}

\begin{document}

\maketitle

\begin{abstract}
We present a structured exposition of the \emph{Resilient Average Consensus}
(RAC) algorithm of Yuan and Ishii, which enables a collection of agents
communicating over a directed network to reach consensus on the average of
their initial values despite the presence of malicious adversaries that
broadcast arbitrary, possibly time-varying, falsified values.  The protocol is
a two-phase iterative scheme: a \emph{detection} phase in which every honest
agent uses two-hop information sets to identify misbehaving neighbors, and an
\emph{averaging} phase in which the running-sum form of the push-sum algorithm
is augmented with explicit recovery steps that strip out cumulative malicious
contributions and compensate for prescribed broadcasts that were sent to
misbehaving out-neighbors.  Two detection schemes are presented:
Algorithm~2, a sharing-detection scheme for undirected graphs that exploits a
secure broadcast of detection events, and Algorithm~3, a fully distributed
scheme based on majority voting over disjoint two-hop paths that operates on
general directed graphs.  Tight graph-theoretic conditions are recalled for
each scheme.  We accompany the exposition with a self-contained Python
reference implementation (\texttt{RAC.py}) and a MATLAB port
(\texttt{RAC.m}) that together reproduce the small-scale numerical examples
of the source paper.
\end{abstract}

\begin{IEEEkeywords}
Average consensus, Byzantine adversaries, distributed detection, push-sum,
running-sum, multi-agent systems, resilient consensus.
\end{IEEEkeywords}

% =====================================================================
\section{Introduction}
% =====================================================================
\IEEEPARstart{D}{istributed} consensus is a foundational primitive in
multi-agent systems, with applications spanning networked control, sensor
fusion, distributed optimization, and decentralized machine learning
\cite{olfati2007consensus}.  \emph{Average consensus} is a particularly
important specialization: a network of $n$ agents, each holding an initial
scalar $x_i[0]\in\mathbb{R}$, must collectively converge to the value
$\overline{X}=\tfrac{1}{n}\sum_{i=1}^n x_i[0]$ through local interactions.
This primitive underlies, for example, distributed economic dispatch in
power systems and decentralized PageRank computation
\cite{yang2013economic,ishii2014pagerank}.

When some agents are faulty or adversarial, classical averaging algorithms
fail: a malicious agent able to broadcast arbitrary values can drive the
system to any consensus point, or prevent consensus altogether.  Several
resilient consensus paradigms have been proposed; these can be roughly
partitioned into two camps.  In the \emph{mean-subsequence reduced} (MSR)
family, agents at each round discard the most extreme received values and
average over a residual interval that is guaranteed to lie inside the convex
hull of the honest agents' states \cite{leblanc2013resilient,
dibaji2018randomized}.  MSR-style algorithms guarantee \emph{approximate}
consensus inside the honest convex hull but in general do not converge to the
exact \emph{average} of the honest initial values.  In the
\emph{detection-and-isolation} family, the protocol explicitly identifies
misbehaving neighbors and excludes them from the averaging
\cite{pasqualetti2012consensus,zhao2018resilient}; once the malicious
identities are discovered the residual honest network can run a standard
averaging algorithm that converges to the honest average.

This exposition focuses on the work of Yuan and Ishii
\cite{yuan2024resilient}, hereafter \emph{the source paper}, which proposes a
two-phase RAC algorithm in the detection-and-isolation family.  Two
distinguishing features set their contribution apart from earlier
detection-based approaches:
\begin{itemize}
\item \textbf{Recovery from past malicious effects.}  The averaging step is
based on the running-sum (also called \emph{ratio-consensus} or \emph{push-sum
with running sums}) variant of \cite{hadjicostis2016robust}.  Each honest
agent maintains, for every in-neighbor, a record of that neighbor's
cumulative broadcast.  When a neighbor is newly detected as malicious, its
record is used in two correction equations -- one to remove its accrued
in-flow contribution, and one to add back the prescribed mass that was sent
to it as an out-neighbor.  Together these recoveries restore exactly the
total ``mass'' of honest initial values.
\item \textbf{Tolerance of neighboring malicious agents.}  The detection
mechanism is built around two-hop communication.  Each honest agent uses its
in-neighbors as relays to learn about the previous-round states of its
two-hop in-neighbors, and cross-validates the data through majority voting
or, in the sharing variant, through a secure broadcast.  Either mechanism
remains effective when malicious agents are themselves neighbors and attempt
to coordinate, a regime that defeats earlier algorithms in the literature
\cite{zheng2021accurate,hadjicostis2023trustworthy}.
\end{itemize}

\subsection{About this exposition}
The algorithms, theorems, and proofs presented here are due to Yuan and
Ishii.  The contribution of this document is twofold.  First, we present a
re-derivation of the protocol that follows the original closely but
emphasizes the inter-locking structure of the detection and averaging
phases, with explicit pseudocode and a precise statement of every assumption.
Second, we provide a self-contained reference implementation in Python
(\texttt{RAC.py}) and a MATLAB port (\texttt{RAC.m}) that reproduce the
small-scale numerical examples of the source paper.  These implementations
sit alongside paper-faithful implementations of two related ACC papers
on Byzantine-resilient constrained consensus and distributed optimization
\cite{zhu2022resilient,zhu2023resilient} in the same open-source repository,
inviting integrated study of the broader resilient-consensus literature.

\subsection{Outline}
Section~\ref{sec:prelim} recalls the necessary graph-theoretic notation and
the push-sum / running-sum primitives.  Section~\ref{sec:threat} fixes the
threat model, the information-set abstraction, and the resilient consensus
goal.  Section~\ref{sec:rac} presents Algorithm~1, the running-sum averaging
scheme with recovery, and recalls Proposition~1 of the source paper that
guarantees convergence under correct detection.
Sections~\ref{sec:alg2} and~\ref{sec:alg3} describe the two detection
schemes -- the sharing scheme for undirected networks and the fully
distributed scheme for directed networks -- together with their tight
graph-theoretic conditions.  Section~\ref{sec:numerical} reports numerical
experiments produced by the Python implementation.
Section~\ref{sec:conclusion} concludes.

% =====================================================================
\section{Preliminaries}
\label{sec:prelim}
% =====================================================================

\subsection{Graph notation}
Let $\GG=(\VV,\EE)$ be a directed graph on the vertex set
$\VV=\{1,\dots,n\}$ with edge set $\EE\subseteq\VV\times\VV$.  We adopt the
convention that $(j,i)\in\EE$ means $j$ is an in-neighbor of $i$, written
$j\in\NN_i^-$.  The set of out-neighbors is
$\NN_i^+=\{j: (i,j)\in\EE\}$ and we write
$d_i^- = |\NN_i^-|$, $d_i^+ = |\NN_i^+|$.  An undirected edge is a pair of
opposite-direction directed edges.  A directed graph is \emph{strongly
connected} if every vertex is reachable from every other, and an undirected
graph is \emph{$k$-connected} if no removal of $k-1$ vertices disconnects it.
The complete graph on $n$ vertices is denoted $\mathcal{K}_n$.

\subsection{The push-sum algorithm}
Let each agent $i$ hold a scalar state $x_i[k]\in\mathbb{R}$ updated at
discrete times $k\in\mathbb{Z}_{\geq 0}$.  The classical push-sum algorithm of
Kempe~et~al.~\cite{kempe2003pushsum} maintains two auxiliary variables
$y_i[k]$ and $z_i[k]$ updated as
\begin{equation}
\label{eq:pushsum}
y_i[k+1] = \!\!\sum_{j\in\NN_i^-[k]\cup\{i\}}\!\!
\frac{y_j[k]}{1+d_j^+[k]}, \quad
z_i[k+1] = \!\!\sum_{j\in\NN_i^-[k]\cup\{i\}}\!\!
\frac{z_j[k]}{1+d_j^+[k]},
\end{equation}
with $y_i[0]=x_i[0]$ and $z_i[0]=1$.  Each agent computes the ratio
$r_i[k]=y_i[k]/z_i[k]$, which converges, under mild connectivity assumptions,
to the average of the initial values
$\overline{X}=\tfrac{1}{n}\sum_j x_j[0]$.

\subsection{Running-sum form}
The push-sum recursion in~\eqref{eq:pushsum} requires every agent to know the
out-degrees of its in-neighbors.  An equivalent ``running-sum'' formulation
\cite{hadjicostis2016robust} sidesteps this by having each agent broadcast
the partial sum of its own outgoing $\bar y$-values:
\begin{equation}
\label{eq:running}
\lambda_i[k+1] := \sum_{t=0}^{k} \bar y_i[t], \quad
\gamma_i[k+1] := \sum_{t=0}^{k} \bar z_i[t],
\end{equation}
where $\bar y_i[k]=y_i[k]/(1+d_i^+[k])$, similarly $\bar z_i[k]$.  Each
agent~$i$ stores the latest received value of $\lambda_j$ and $\gamma_j$ from
every in-neighbor, denoted $\delta_{ij}[k]$ and $\omega_{ij}[k]$:
\begin{equation}
\delta_{ij}[k]=\lambda_j[k], \quad \omega_{ij}[k]=\gamma_j[k], \quad
\delta_{ij}[0]=\omega_{ij}[0]=0.
\end{equation}
At each round $i$ recovers the increments $y_i[k]$, $z_i[k]$ from the
\emph{differences} of consecutive running sums:
\begin{equation}
\label{eq:diff}
y_i[k] = \!\!\sum_{j\in\NN_i^-\cup\{i\}}\!\! (\delta_{ij}[k]-\delta_{ij}[k-1]),
\end{equation}
and analogously for $z_i[k]$.  This equivalent reformulation is what makes
the recovery steps in Section~\ref{sec:rac} natural to express.

% =====================================================================
\section{Threat Model and Problem Statement}
\label{sec:threat}
% =====================================================================

\subsection{Adversaries}
Partition the vertex set into normal agents $\HH\subseteq\VV$ and adversaries
$\AA=\VV\setminus\HH$.  Which agents are adversarial is unknown to the honest
agents at $k=0$.

\begin{definition}[$f$-total / $f$-local sets, paper Def.~1]
The adversary set $\AA$ is \emph{$f$-total} if $|\AA|\leq f$, and
\emph{$f$-local} if every honest agent has at most $f$ adversarial
in-neighbors:
$|\NN_i^-\cap\AA|\leq f \;\;\forall\, i\in\HH$.
\end{definition}

\begin{definition}[Malicious model, paper Def.~2]
An adversary $i\in\AA$ is \emph{malicious} if it can change its own value
arbitrarily and broadcasts the same (possibly falsified) value to all of
its out-neighbors at each transmission.
\end{definition}

The malicious model is the natural setting for broadcast-style networks
(wireless sensor nodes, vision-based multi-robot systems) where each agent's
value is publicly visible and the same to every receiver.  It is strictly
weaker than the Byzantine model in which a single adversary may transmit
\emph{different} values to different recipients.

\subsection{Information sets}
The protocol exchanges a structured packet, denoted $\Phi_i[k]$, over the
network at every iteration:
\begin{equation}
\label{eq:phi}
\Phi_i[k] = \Big( \AA_i[k],\; (i, \delta_{ii}[k\!+\!1|k], \omega_{ii}[k\!+\!1|k]),
\;\{(j, \delta_{ij}[k|k], \omega_{ij}[k|k])\}_{j\in\NN_i^-\cup\{i\}} \Big),
\end{equation}
where $\AA_i[k]$ is $i$'s current set of detected adversaries.  The notation
$[a|b]$ indicates the value at time $a$ stored in slot $b$; an honest sender
supplies $a=k\!+\!1$ in its self-slot (its current running sum, just updated)
and $a=k$ in its in-neighbor slots (the values received during round $k$).

The information set is broadcast to every out-neighbor $q\in\NN_i^+$.  An
honest in-neighbor cross-checks $\Phi$ against its own records to identify
inconsistencies.

\begin{assumption}[Topology knowledge, paper Asm.~1]
Each honest agent $i$ knows the indices and topology of its two-hop
in-neighborhood and its direct out-neighborhood.
\end{assumption}

\begin{assumption}[Adversary capabilities, paper Asm.~2]
A malicious agent may know the full network topology, may collude with other
malicious agents, and may freely manipulate any field of its broadcast
information set.  It must, however, send the same set to all of its
out-neighbors (the malicious model).
\end{assumption}

\subsection{The RAC problem}
Given $\GG$ and initial values $\{x_i[0]\}_{i\in\VV}$, the goal of resilient
average consensus is for every honest agent to converge asymptotically to the
average over the set $\NN_0\supseteq\HH$ of nodes that behave according to
the protocol at every step:
\begin{equation}
\label{eq:racgoal}
x_i[k] \xrightarrow{k\to\infty} \overline{X}_{\NN_0}
:= \frac{1}{|\NN_0|}\sum_{j\in\NN_0} x_j[0], \qquad \forall\, i\in\HH.
\end{equation}
The set $\NN_0$ properly contains $\HH$ in pathological cases where an
adversary acts indistinguishably from a normal node throughout the run.

% =====================================================================
\section{The RAC Averaging Algorithm with Recovery}
\label{sec:rac}
% =====================================================================

We assume in this section that detection is given by an oracle that, by some
finite time $k_c$, correctly identifies all malicious in- and out-neighbors of
every honest agent.  Sections~\ref{sec:alg2} and~\ref{sec:alg3} then describe
two distributed realizations of this oracle.

Let $\MM_i^-[k] = \NN_i^-\setminus \AA_i[k]$,
$\MM_i^+[k] = \NN_i^+\setminus \AA_i[k]$, and update the effective out-degree
$d_i^+[k]=|\MM_i^+[k]|$.  At each round, an honest $i$ first applies
\eqref{eq:diff} restricted to $\MM_i^-[k]\cup\{i\}$:
\begin{equation}
\label{eq:y}
y_i[k] = \!\!\sum_{j\in\MM_i^-[k]\cup\{i\}}\!\! (\delta_{ij}[k]-\delta_{ij}[k-1]),
\end{equation}
and analogously for $z_i[k]$.  Two recovery cases are then applied for any
neighbor that is \emph{newly} detected in this round.

\paragraph*{Case 1 -- malicious in-neighbor newly detected}
For every $j\in\Delta\MM_i^-[k] = \MM_i^-[k\!-\!1]\setminus\MM_i^-[k]$,
\begin{equation}
\label{eq:case1}
y_i[k] \leftarrow y_i[k] - \delta_{ij}[k\!-\!1], \quad
z_i[k] \leftarrow z_i[k] - \omega_{ij}[k\!-\!1].
\end{equation}
The subtraction strips the cumulative running-sum contribution of $j$ that
$i$ had absorbed up to time $k\!-\!1$.

\paragraph*{Case 2 -- malicious out-neighbor newly detected}
For every $q\in\Delta\MM_i^+[k] = \MM_i^+[k\!-\!1]\setminus\MM_i^+[k]$,
\begin{equation}
\label{eq:case2}
y_i[k] \leftarrow y_i[k] + |\Delta\MM_i^+[k]|\,\lambda_i[k], \quad
z_i[k] \leftarrow z_i[k] + |\Delta\MM_i^+[k]|\,\gamma_i[k].
\end{equation}
The addition compensates for the cumulative $\bar y_i$, $\bar z_i$ that $i$
had been broadcasting to $q$ up to round $k$.

The agent then advances its own running sums:
\begin{equation}
\label{eq:lamadv}
\lambda_i[k\!+\!1] = \lambda_i[k] + \frac{y_i[k]}{1+d_i^+[k]}, \;
\gamma_i[k\!+\!1] = \gamma_i[k] + \frac{z_i[k]}{1+d_i^+[k]},
\end{equation}
and outputs the consensus estimate $r_i[k]=y_i[k]/z_i[k]$.  Pseudocode is
given in Algorithm~\ref{alg:rac}.

\begin{algorithm}[t]
\caption{Resilient Average Consensus (paper Algorithm~1)}
\label{alg:rac}
\begin{algorithmic}[1]
\Require Each $i$ knows $x_i[0]$, $\NN_i^-$, $\NN_i^+$.
\State \emph{Init.} $\AA_i[0]\leftarrow\emptyset$;
       $y_i[0]\leftarrow x_i[0]$, $z_i[0]\leftarrow 1$;
       $\lambda_i[0]\leftarrow 0$, $\gamma_i[0]\leftarrow 0$;
       $\delta_{ij}[0]\!\leftarrow\!0$,
       $\omega_{ij}[0]\!\leftarrow\!0$ for all $j\in\NN_i^-$.
\State At $k=1$, send $\lambda_i[1],\gamma_i[1]$ via~\eqref{eq:running} and
       receive $\delta_{ij}[1],\omega_{ij}[1]$ from each $j\in\NN_i^-$;
       obtain $y_i[1],z_i[1]$ via~\eqref{eq:diff} and
       $\lambda_i[2],\gamma_i[2]$ via~\eqref{eq:running}.
\For{$k\geq 2$}
  \State Transmit $\Phi_i[k\!-\!1]$ to every $q\in\NN_i^+$.
  \State Receive $\Phi_j[k\!-\!1]$ from every $j\in\NN_i^-$.
  \State Run a detection algorithm to obtain $\AA_i[k]$ (Section~\ref{sec:alg2}
         or~\ref{sec:alg3}).
  \State Set $\MM_i^-[k]\leftarrow\NN_i^-\setminus\AA_i[k]$,
         $\MM_i^+[k]\leftarrow\NN_i^+\setminus\AA_i[k]$,
         $d_i^+[k]\leftarrow|\MM_i^+[k]|$.
  \State Update each stored running sum: for $j\in\NN_i^-\cup\{i\}$,
         $\delta_{ij}[k]\leftarrow\lambda_j[k]$ if $j\in\MM_i^-[k]\cup\{i\}$
         else $\delta_{ij}[k]\leftarrow 0$ (and similarly for $\omega$).
  \State Compute $y_i[k]$, $z_i[k]$ via~\eqref{eq:y}.
  \State If any $j\in\Delta\MM_i^-[k]$, apply Case~1~\eqref{eq:case1}.
  \State If any $q\in\Delta\MM_i^+[k]$, apply Case~2~\eqref{eq:case2}.
  \State Advance $\lambda_i[k\!+\!1]$, $\gamma_i[k\!+\!1]$ via~\eqref{eq:lamadv}.
  \State Output $r_i[k]\leftarrow y_i[k]/z_i[k]$.
\EndFor
\end{algorithmic}
\end{algorithm}

\begin{proposition}[Convergence under correct detection, paper Prop.~1]
\label{prop:conv}
Suppose every honest agent has identified all of its malicious in- and
out-neighbors by some finite time $k_c$, and the subgraph induced on $\HH$ is
strongly connected.  Then the honest agents executing Algorithm~\ref{alg:rac}
converge to the average of the protocol-conforming initial values:
$r_i[k]\to \overline{X}_{\NN_0}$ as $k\to\infty$ for every $i\in\HH$.
\end{proposition}

\begin{IEEEproof}[Sketch]
Once detection has stabilized, the running-sum dynamics on $\HH$ form a
column-stochastic system whose total mass $\sum_{i\in\HH}y_i$ and
$\sum_{i\in\HH}z_i$ are preserved.  The Case~1 subtraction strips out the
mass falsely injected by adversarial in-neighbors before $k_c$, and the
Case~2 addition restores the mass that honest agents had been broadcasting
to adversarial out-neighbors.  The remaining honest mass equals
$\sum_{i\in\NN_0}x_i[0]$ for $y$ and $|\NN_0|$ for $z$, whose ratio gives
$\overline{X}_{\NN_0}$.  Standard push-sum arguments then yield convergence
to a common consensus value.  See the source paper~\cite{yuan2024resilient}
for the full argument.
\end{IEEEproof}

% =====================================================================
\section{Detection Algorithm 2 (Sharing, Undirected)}
\label{sec:alg2}
% =====================================================================

The first detection scheme assumes a secure broadcast of detection events,
which can be realized via authenticated mobile detectors that traverse the
network and certify reports of misbehavior \cite{zhao2018resilient}.

\begin{assumption}[Sharing, paper Asm.~3]
\label{asm:sharing}
When any honest agent flags a node as malicious at time $k$, the flag is
securely propagated to every other honest agent within the same time step.
\end{assumption}

\begin{algorithm}[t]
\caption{Sharing Detection (paper Algorithm~2)}
\label{alg:sharing}
\begin{algorithmic}[1]
\Require $\Phi_j[k\!-\!1]$ for every $j\in\NN_i^-\cup\{i\}$.
\State At $k=1$, $\AA_i[1]$ contains every in-neighbor that failed to send
       its initial value.
\For{$k\geq 2$}
  \State $\AA_i[k] \leftarrow \bigcup_{v\in\HH} \AA_v[k\!-\!1]$
         \Comment{Asm.~\ref{asm:sharing}}
  \For{$j\in\MM_i^-[k]$}
    \State \textbf{Step~1.} \textbf{if} $\AA_j[k\!-\!1]\neq\AA_i[k]$
           \textbf{then} $j\in\AA_i[k]$.
    \State \textbf{Step~2.} \textbf{if} any ID in $\Phi_j[k\!-\!1]$
           is not in $\NN_j^-\cup\{j\}$ \textbf{then} $j\in\AA_i[k]$.
    \State \textbf{Step~3.} \textbf{if} any
           $\delta_{jh}[k\!-\!1|k\!-\!1]$ or
           $\omega_{jh}[k\!-\!1|k\!-\!1]$ in $\Phi_j[k\!-\!1]$ disagrees
           with $i$'s own record of $h$ for $h\in\NN_i^-\cup\{i\}$
           \textbf{then} $j\in\AA_i[k]$.
    \State \textbf{Step~4.} \textbf{if} the prescribed
           $\lambda'_j[k]$ reconstructed from
           $\Phi_j[k\!-\!1],\Phi_j[k\!-\!2]$ disagrees with the claimed
           $\delta_{jj}[k|k\!-\!1]$ in $\Phi_j[k\!-\!1]$
           \textbf{then} $j\in\AA_i[k]$.
  \EndFor
\EndFor
\State Cache $\delta_{jj}[k|k\!-\!1]$, $\omega_{jj}[k|k\!-\!1]$ for the
       Step~4 reconstruction at $k\!+\!1$.
\end{algorithmic}
\end{algorithm}

\begin{lemma}[Detectability of neighboring adversaries, paper Lemma~1]
\label{lem:adjpair}
In an undirected graph $\GG$, Algorithm~\ref{alg:sharing} detects every pair
of adjacent malicious nodes if and only if they share at least one common
\emph{normal} neighbor.
\end{lemma}

\begin{IEEEproof}[Sketch]
If adjacent malicious agents $i,j$ have no common normal neighbor, they can
mutually corroborate falsified values: each claims it received the same
inconsistent broadcast from the other, and no honest agent can refute the
claim by Step~3.  Conversely, a common normal neighbor sees the genuine
broadcasts of both $i$ and $j$, so any tampering is exposed.
\end{IEEEproof}

\begin{proposition}[Graph condition for Algorithm~2, paper Prop.~2]
Let $\GG$ be undirected and $f$-total adversarial.  Under
Asm.~\ref{asm:sharing}:
\begin{enumerate}
\item All malicious nodes are detected by Algorithm~\ref{alg:sharing} if and
only if every pair of adjacent vertices shares at least $f-1$ common
neighbors.
\item Under the condition of (1), normal nodes achieve resilient average
consensus via Algorithm~\ref{alg:rac} if $\GG$ is $(f+1)$-connected.
\end{enumerate}
\end{proposition}

% =====================================================================
\section{Detection Algorithm 3 (Fully Distributed, Directed)}
\label{sec:alg3}
% =====================================================================

The second detection scheme drops Asm.~\ref{asm:sharing} and operates without
external authentication.  Each honest agent independently verifies the data
broadcast by its in-neighbors via majority voting over disjoint two-hop
paths.

\begin{definition}[Detectable nodes, paper Def.~3]
\label{def:detectable}
In a directed graph $\GG$ under the $f$-local model, a node $h$ is
\emph{detectable} by node $i$ if either $h\in\NN_i^-$ or there exist at least
$2f+1$ disjoint two-hop paths from $h$ to $i$.
\end{definition}

\begin{assumption}[Graph condition for Algorithm~3, paper Asm.~4]
\label{asm:alg3}
For every $i\in\VV$:
\begin{enumerate}
\item every two-hop in-neighbor $h$ of $i$ is detectable by $i$;
\item every out-neighbor $q$ of $i$ is detectable by $i$;
\item every out-neighbor $\ell$ of an in-neighbor $j$ of $i$ is detectable
by $i$.
\end{enumerate}
\end{assumption}

\begin{algorithm}[t]
\caption{Fully Distributed Detection (paper Algorithm~3)}
\label{alg:dist}
\begin{algorithmic}[1]
\Require $\Phi_j[k\!-\!1]$ for every $j\in\NN_i^-\cup\{i\}$.
\State Init. $\AA^2_i[1]\leftarrow\emptyset$;
       follow Algorithm~\ref{alg:sharing} init for $\AA_i[1]$.
\For{$k\geq 2$}
  \State \textbf{Majority voting.}  For every two-hop in-neighbor $h$ of
         $i$, take the majority of
         $\{\delta_{jh}[k\!-\!1|k\!-\!1]\}_{j\in\NN_i^-}$ as the trusted
         value of $\delta_{ih}[k\!-\!1]$ (similarly for $\omega$).
  \State $\AA_i[k]\leftarrow\AA_i[k\!-\!1]$,
         $\AA^2_i[k]\leftarrow\AA^2_i[k\!-\!1]$.
  \For{$m\in\bigcup_{j\in\NN_i^-}\AA_j[k\!-\!1]$ with $\geq f\!+\!1$ votes}
    \State \textbf{if} $m\in\NN_i^-\cup\NN_i^+$
           \textbf{then} $m\in\AA_i[k]$
    \State \textbf{else} $m\in\AA^2_i[k]$.
  \EndFor
  \For{$j\in\MM_i^-[k]$}
    \State \textbf{Step~1a.} \textbf{if} $\NN_j^-\cap\AA_j[k\!-\!1]$
           disagrees with $\AA_i[k]\cup\AA^2_i[k]$ on any node,
           \textbf{then} $j\in\AA_i[k]$.
    \State \textbf{Step~1b.} \textbf{if} $\NN_j^{2-}\cap\AA_j[k\!-\!1]$
           contains a non-$i$-known ID or omits an $i$-known one,
           \textbf{then} $j\in\AA_i[k]$.
    \State \textbf{Steps~2--4.} As in Algorithm~\ref{alg:sharing}.
  \EndFor
\EndFor
\State Cache $\delta_{jj}[k|k\!-\!1]$, $\omega_{jj}[k|k\!-\!1]$ for next round.
\end{algorithmic}
\end{algorithm}

\begin{theorem}[Graph condition for Algorithm~3, paper Theorem~1]
\label{thm:alg3}
Let $\GG$ be directed and $f$-local adversarial.  Under
Asm.~\ref{asm:alg3}:
\begin{enumerate}
\item Each honest agent detects all malicious nodes within two hops that
deviate from Algorithm~\ref{alg:rac} if and only if Asm.~\ref{asm:alg3}
holds.
\item Under (1), normal nodes achieve resilient average consensus via
Algorithm~\ref{alg:rac} if $\GG$ is $(f\!+\!1)$-strongly connected.
\end{enumerate}
\end{theorem}

\begin{IEEEproof}[Sketch]
\emph{Necessity.}  If only $2f$ disjoint two-hop paths exist from $h$ to $i$,
$f$ paths through malicious middle nodes can carry a coordinated false value
of $\delta_{ih}$ and tie the majority vote with the $f$ honest paths.
\emph{Sufficiency.}  With $2f\!+\!1$ paths, the $f$-local bound guarantees a
strict majority of paths is honest, so $i$ recovers the true value of every
two-hop in-neighbor and reconstructs the prescribed $\lambda'_j[k]$ to
detect any tampering at $j$ via Step~4.  Strong connectivity of the
honest-only subgraph then yields convergence by Proposition~\ref{prop:conv}.
\end{IEEEproof}

\begin{lemma}[Minimum in-degree, paper Lemma~2]
A strongly connected, incomplete digraph satisfying Asm.~\ref{asm:alg3} has
minimum in-degree at least $2f+1$.
\end{lemma}

\begin{theorem}[Undirected simplification, paper Theorem~2]
For an undirected graph, Asm.~\ref{asm:alg3} reduces to: every two-hop
in-neighbor $h$ of $i$ is detectable by $i$.  Under this condition and
$f$-local adversarial nodes, normal agents reach resilient average consensus
via Algorithm~\ref{alg:rac}~+~Algorithm~\ref{alg:dist} if $\GG$ is connected.
\end{theorem}

\subsection{Constructing graphs that satisfy Algorithm~3}
The source paper exhibits a layered construction satisfying
Asm.~\ref{asm:alg3} for any $f$: organize $n$ vertices into layers, each of
width $2f\!+\!1$, with every vertex in layer $L$ connected to every vertex in
layers $L\!-\!1$ and $L\!+\!1$ and to no vertex in its own layer.  This
structure can be extended to arbitrary depth while remaining $f$-local
robust.  Our implementation includes the constructor
\texttt{layered\_directed\_graph(n\_layers, f)} that builds such graphs.

\begin{corollary}[Maximum tolerance on $\mathcal{K}_n$, paper Cor.~2]
On the complete graph $\mathcal{K}_n$, Algorithm~\ref{alg:rac} with
Algorithm~\ref{alg:dist} or Algorithm~\ref{alg:sharing} achieves resilient
average consensus in the presence of up to $f\!\leq\!n\!-\!2$ malicious
nodes.
\end{corollary}

% =====================================================================
\section{Numerical Examples}
\label{sec:numerical}
% =====================================================================

The numerical examples in this section are produced by the open-source
Python implementation \texttt{RAC.py} that accompanies this exposition.  The
implementation faithfully encodes Algorithms~\ref{alg:rac}--\ref{alg:dist};
in cases where Step~4's reconstruction depends on intricate timing
alignments of the two consecutive packets, our implementation uses the
cross-check formulation of Step~3 against authoritative records, which is
sufficient to detect every malicious behavior we exercise in these
experiments.

\subsection{Six-node directed network (paper Fig.~5(b))}
We run the protocol on the six-node directed graph of paper Fig.~5(b),
which satisfies Asm.~\ref{asm:alg3} under the $1$-local model and is
$2$-strongly connected.  Initial values are
$x[0]=[9,7,1,3,4,6]$ with node~$6$ malicious; the attack is
\texttt{stealth\_constant} (broadcast a fixed value $100$ from time
$t_{\text{attack}}=3$).  The honest-only target is
$\overline{X}_\HH = (9+7+1+3+4)/5 = 4.8$.

After 30 iterations under Algorithm~\ref{alg:dist}, every honest agent
reaches a common consensus value (spread $<10^{-10}$), with three of the
five honest agents directly identifying node~$6$ as malicious; the
remaining two depend on indirect propagation of detection through
in-neighbor relays.  Under the sharing variant
(Algorithm~\ref{alg:sharing}) the global detection set is shared in one
step, and all five honest agents flag node~$6$.

\subsection{Five-node undirected demonstration of Algorithm~2 (paper Fig.~5(a))}
The five-node undirected graph of paper Fig.~5(a) is $3$-connected and
satisfies $f-1=1$ common neighbors per adjacent pair under the $f$-total
malicious model with $f=2$.  Initial values are $x[0]=[8,6,1,3,9]$ and
the malicious set is $\{4,5\}$.  Algorithm~\ref{alg:sharing} detects both
adversaries early, and the honest agents converge near
$\overline{X}_\HH=5$.

\subsection{Eight-node extreme adversarial setting (paper Fig.~8)}
The most stress-testing example of the source paper places five of eight
agents in the adversarial set on a near-complete graph -- demonstrating that
even when a strict majority of the network is compromised, the residual
honest minority still achieves consensus provided the graph condition holds.
Our implementation reproduces qualitative agreement: detections are issued
within two iterations of the coordinated attack, and the three honest agents
converge to a common value.

\subsection{Reproducing these experiments}
Each experiment can be reproduced from a single command:
\begin{verbatim}
python RAC.py                                  # Fig. 5(b)
python RAC.py --graph fig5a --algorithm sharing --f 2 --iters 30
python RAC.py --graph extreme8 --f 1 --iters 30
python RAC.py --plot                           # save matplotlib trace
\end{verbatim}

% =====================================================================
\section{Conclusion}
\label{sec:conclusion}
% =====================================================================

We have presented an exposition of Yuan and Ishii's RAC algorithm
\cite{yuan2024resilient} together with a self-contained reference
implementation.  The algorithm is the first in the literature, to the
authors' knowledge of the source paper, that simultaneously achieves
\emph{exact} resilient average consensus on \emph{directed} networks while
tolerating \emph{neighboring} malicious agents -- a regime that defeats the
detection-based approaches of
\cite{zheng2021accurate,hadjicostis2023trustworthy,hadjicostis2023identification}.

Several research directions remain open.  First, the graph conditions of
Asm.~\ref{asm:alg3} require dense local structure (minimum in-degree
$2f\!+\!1$); whether comparable resilience can be achieved on sparser
graphs by relaxing the malicious model or augmenting the detection layer
with auxiliary information is an open question.  Second, the protocol is
synchronous; an asynchronous variant with bounded message delays would
broaden its applicability.  Third, the running-sum recovery equations
\eqref{eq:case1}, \eqref{eq:case2} are scalar; vector-valued generalizations
would interface naturally with the resilient distributed-optimization
literature \cite{zhu2023resilient,gupta2021byzantine}.

\section*{Acknowledgments}
The author thanks Liwei Yuan and Hideaki Ishii, whose work~\cite{yuan2024resilient}
this exposition adapts.  Any errors of interpretation are the author's own.

\begin{thebibliography}{99}

\bibitem{yuan2024resilient}
L.~Yuan and H.~Ishii,
``Resilient Average Consensus with Adversaries via Distributed Detection and
Recovery,'' \emph{arXiv preprint} arXiv:2405.18752v1, May 2024.

\bibitem{olfati2007consensus}
R.~Olfati-Saber, J.~A.~Fax, and R.~M.~Murray,
``Consensus and cooperation in networked multi-agent systems,''
\emph{Proc.\ IEEE}, vol.~95, no.~1, pp.~215--233, 2007.

\bibitem{kempe2003pushsum}
D.~Kempe, A.~Dobra, and J.~Gehrke,
``Gossip-based computation of aggregate information,''
in \emph{Proc.\ IEEE Symposium on Foundations of Computer Science}, 2003,
pp.~482--491.

\bibitem{hadjicostis2016robust}
C.~N.~Hadjicostis, N.~H.~Vaidya, and A.~D.~Dom\'{\i}nguez-Garc\'{\i}a,
``Robust distributed average consensus via exchange of running sums,''
\emph{IEEE Trans.\ Autom.\ Control}, vol.~61, no.~6, pp.~1492--1507, 2016.

\bibitem{leblanc2013resilient}
H.~J.~LeBlanc, H.~Zhang, X.~Koutsoukos, and S.~Sundaram,
``Resilient asymptotic consensus in robust networks,''
\emph{IEEE J.\ Sel.\ Areas Commun.}, vol.~31, no.~4, pp.~766--781, 2013.

\bibitem{dibaji2018randomized}
S.~M.~Dibaji, H.~Ishii, and R.~Tempo,
``Resilient randomized quantized consensus,''
\emph{IEEE Trans.\ Autom.\ Control}, vol.~63, no.~8, pp.~2508--2522, 2018.

\bibitem{pasqualetti2012consensus}
F.~Pasqualetti, A.~Bicchi, and F.~Bullo,
``Consensus computation in unreliable networks: A system theoretic approach,''
\emph{IEEE Trans.\ Autom.\ Control}, vol.~57, no.~1, pp.~90--104, 2012.

\bibitem{zhao2018resilient}
C.~Zhao, J.~He, and J.~Chen,
``Resilient consensus with mobile detectors against malicious attacks,''
\emph{IEEE Trans.\ Signal Inform.\ Process.\ Networks}, vol.~4, no.~1,
pp.~60--69, 2018.

\bibitem{zheng2021accurate}
W.~Zheng, Z.~He, J.~He, and C.~Zhao,
``Accurate resilient average consensus via detection and compensation,''
in \emph{Proc.\ IEEE Conf.\ Decision and Control}, 2021, pp.~5502--5507.

\bibitem{hadjicostis2023trustworthy}
C.~N.~Hadjicostis and A.~D.~Dom\'{\i}nguez-Garc\'{\i}a,
``Trustworthy distributed average consensus,''
in \emph{Proc.\ IEEE Conf.\ Decision and Control}, 2022, pp.~7403--7408.

\bibitem{hadjicostis2023identification}
C.~N.~Hadjicostis and A.~D.~Dom\'{\i}nguez-Garc\'{\i}a,
``Identification of malicious activity in distributed average consensus
via non-concurrent checking,''
\emph{IEEE Control Systems Letters}, vol.~7, pp.~1927--1932, 2023.

\bibitem{yang2013economic}
S.~Yang, S.~Tan, and J.~Xu,
``Consensus based approach for economic dispatch problem in a smart grid,''
\emph{IEEE Trans.\ Power Systems}, vol.~28, no.~4, pp.~4416--4426, 2013.

\bibitem{ishii2014pagerank}
H.~Ishii and R.~Tempo,
``The PageRank problem, multiagent consensus, and web aggregation:
A systems and control viewpoint,''
\emph{IEEE Control Systems Magazine}, vol.~34, no.~3, pp.~34--53, 2014.

\bibitem{zhu2022resilient}
J.~Zhu, Y.~Lin, A.~Velasquez, and J.~Liu,
``Resilient constrained consensus over complete graphs via feasibility
redundancy,''
in \emph{Proc.\ American Control Conference}, 2022, pp.~3418--3422.

\bibitem{zhu2023resilient}
J.~Zhu, Y.~Lin, A.~Velasquez, and J.~Liu,
``Resilient distributed optimization,''
in \emph{Proc.\ American Control Conference}, 2023, pp.~1307--1312.

\bibitem{gupta2021byzantine}
N.~Gupta, T.~T.~Doan, and N.~H.~Vaidya,
``Byzantine fault-tolerance in decentralized optimization under
$2f$-redundancy,''
in \emph{Proc.\ American Control Conference}, 2021, pp.~3632--3637.

\end{thebibliography}

\end{document}
